\chapter{Defeasible reasoning and belief revision in psychology}
\label{sec:p19-rationality-5-4}

Gazzo Castañeda and Knauff's 2021 chapter distinguishes withdrawing a
conclusion from changing the premises used to derive it. It then asks
whether that distinction identifies different cognitive processes or partly
reflects the tasks used to investigate them. Its central proposal is that
retrieving a disabling circumstance can contribute to both activities.
The review compares mental-model, probabilistic and dual-process explanations
without supplying a single experimentally settled mechanism
\cite[pp.~347--359]{GazzoCastanedaKnauff2021Defeasible}.

This treatment follows that source's full argument, examples and note.
The detailed primary records below concern Wolf and Knauff's two 2008
experiments and Knauff and colleagues' 2013 Experiment~1. Byrne's
Experiment~1 has its own record in
\cref{sec:p19-conditional-reasoning-suppression}; Cariani and Rips's later
interpretation experiments remain separate studies in
\cref{sec:p19-conditional-interpretation-experiments}.
These records are not pooled into an estimate of a general human effect.
The broader findings discussed at review level below do not acquire
participant denominators or independently checked uncertainty estimates
merely by appearing in the chapter's bibliography.

\section{Two changes, two laboratory tasks}
\label{sec:p19-r54-two-tasks}

\begin{hbdefinition}{def:p19-r54-defeat-revision}{Derived conclusions, prior beliefs and imposed evidence}
In the chapter's usage, defeasible reasoning retracts a previously drawn
conclusion when additional circumstances undermine it, while retaining a
more general prior belief. Belief revision instead changes prior beliefs
to accommodate conflicting evidence. In the common laboratory revision
task, the premise sentences stand in for prior beliefs; an additional
conflicting statement is declared certain, and participants indicate which
earlier premise they believe less. This finite sentence-choice task is a
strong simplification of a person's whole belief system.

The examples concern a person expecting tiredness after a late bedtime
but learning that sleeping in is possible, and a person expecting meditation
to aid sleep but remaining unable to sleep. The first suggests a disabling
circumstance; the second can prompt doubt about the advice, its application,
or another background assumption. The examples illustrate the distinction,
not experiments with measured samples.
\end{hbdefinition}

For fixed material premises, MP and MT are valid; AC and DA require additional
assumptions. That benchmark does not determine how an English conditional
is interpreted. Nor does it decide whether a response measures necessity,
confidence, relative premise credibility or an explanation of an exception.
In particular, a button indicating the more credible premise is not a
direct observation that the other sentence was removed from a deductively
closed theory.

\begin{hblemma}{thm:p19-r54-disabler-boundary}{proved=annex:p19.r54.disabler-boundary}{Adding an explanation is not a consistency repair of unchanged material beliefs}
The material base \(\{p\to q,p,\neg q\}\) is inconsistent.
Adding a disabling conditional \(r\to\neg q\), or adding r,
does not make it consistent. By contrast
\(\{(p\land\neg r)\to q,p,\neg q,r\}\) is consistent.
Thus the source's explanation by a wrong kind of meditation must revise
an assumption, reinterpret a conditional as defeasible, or change another
represented commitment; it is not an unchanged material-premise deduction.
The same distinction applies to keeping a general default about bedtime
while withdrawing its conclusion for someone who can sleep in.
\end{hblemma}

\section{What the suppression paradigms measure}
\label{sec:p19-r54-paradigms}

In the \emph{overt} paradigm, a simple conditional/categorical argument
is compared with an argument containing explicit extra information. A
disabler makes the antecedent less sufficient for the consequence;
an alternative supplies another route to the consequence and makes the
antecedent less necessary. The chapter's bedtime examples contrast
sleeping in with another reason for tiredness, such as disrupted sleep.
The hypothesized selective contrasts are MP/MT for disablers and AC/DA
for alternatives. Added premises are not assumed interchangeable simply
because the surface strings have the same grammatical pattern
\cite[\S2.1]{GazzoCastanedaKnauff2021Defeasible}.
In the source alternative example, not going to bed late does not ensure
not being tired when disrupted sleep is another route to tiredness. This
is the DA target, not the MP target of the sleeping-in example.

In the \emph{covert} paradigm, one group generates possible disablers and
alternatives for conditionals. Materials are categorized by many/few
generated disablers and many/few alternatives. A different group assesses
inferences from these conditionals without being shown those generated
cases. The second group's responses can therefore depend on accessible
background information, but a material's norming-group count is not that
individual participant's measured retrieval count. Between-material
associations do not by themselves establish which representation or
memory operation generated each response
\cite[\S2.2]{GazzoCastanedaKnauff2021Defeasible}.
The generation prompt accordingly contrasts a late bedtime without
tiredness, to elicit disablers, with tiredness without a late bedtime,
to elicit alternatives. The two prompts seek different kinds of case.

The review organizes the content evidence around familiarity, expertise,
specificity and associative strength. It distinguishes the number of
\emph{different} exception types from the frequency of exceptions: a
refrigerator light can fail for several reasons even if failures are rare.
Its discussion contrasts the gradual effects attributed to De Neys and
colleagues with Geiger and Oberauer's frequency account. The chapter's
single note limits the latter evidence to \emph{disablers}; an analogous
frequency effect for alternatives is suggested, not established by that
study \cite[\S2.3 and note~1]{GazzoCastanedaKnauff2021Defeasible}.
Neither an inverse relation between endorsement and exception count nor
its description as ``monotonic'' means monotonic logical consequence.
The source's association example contrasts sleeping in with a cold
morning shower as different potential ways to avert tiredness. Its
refrigerator example distinguishes broken equipment, a broken bulb and
absent electricity as exception types from how often an exception occurs.
The number of named causes and their estimated frequency are different
predictor quantities, not two labels for one measured variable.

Instructions, response scales and expressions of credibility form a
separate context dimension. A graded confidence judgment can express
differences that a binary answer cannot. A dichotomous task can also induce
a different response criterion or strategy. The review cites work on
necessity instructions, response modality, conditional certainty and
confidence in a disabler; these are not one controlled manipulation
pooled across interchangeable populations
\cite[\S2.4]{GazzoCastanedaKnauff2021Defeasible}.
The review's fishing example adds the requirement of catching a fish
before a fish supper, then varies confidence in that disabling possibility
by describing the fisher as lucky. Its uncertainty-language example
inserts high- or low-probability wording into a conditional. These
illustrations show what the linguistic manipulation concerns; they do
not supply a new sample or an effect size in this chapter.
The later Cariani--Rips studies already recorded in P19 are a useful
contrast: suppression also occurred under explicitly deductive instructions
in their particular samples. This does not prove that instructions never
matter, and their different premises and response targets remain visible.

\section{Revision, credibility and the spatial contrast}
\label{sec:p19-r54-revision-evidence}

For an MP revision problem the review presents an uncertain conditional,
its uncertain antecedent, and a newly certain negated consequent. For MT
it presents the conditional, a negated consequent, and a newly certain
antecedent. The response asks which uncertain premise to abandon or
believe less. The certainty of the new input is a task instruction, not
an independently demonstrated psychological commitment.

The chapter contrasts early preferences for retaining general statements
with later conditional tasks in which the general conditional was more
often questioned. Its proposed reconciliation appeals to familiar,
plausible content versus unfamiliar material. It also distinguishes
source trustworthiness and the MP--MT asymmetry, arguing that content
alone does not explain every response contrast. These are the review's
claim-level comparisons among different studies, not a common response
variable measured in a single sampled population
\cite[\S3.1]{GazzoCastanedaKnauff2021Defeasible}.
The source illustrates a familiar conditional with water extinguishing
a campfire, and an unfamiliar one with fictitious Pings and forcefields.
Its earlier quantified examples concern a general vertebrate/backbone
statement and a particular biological classification. These mark the
content contrast in the review; they are not silently assumed identical
to every administered stimulus of the cited experiments. Retaining a
familiar generalization need not be the same task as retaining an
unfamiliar if--then sentence.

\subsection{Wolf and Knauff 2008: actual samples and incomplete second experiment}

The conference paper's Experiment~1 tested 30 Giessen students, 17 women
and 13 men, aged 19--31, for monetary compensation. Thirty-six conditionals
had first been rated by another group from the same student population;
the paper does not state that group's size. Eight low-rated and eight
high-rated statements were selected, with reported average ratings 7.19
and 88.44 on the 0--100 scale. This was selection of content by a pretest,
not assignment of numerical probabilities 0 and 100 to identical events.

Each participant answered 32 problem sets in a within-person
MP/MT by high/low design. Three sentences appeared sequentially; the
first two were declared uncertain and the conflicting third certainly
true. Participants chose which of the first two they believed more.
Item presentation was randomized and response-button assignment
counterbalanced. The measures were relative-premise choice and decision
time \cite[pp.~831--834]{WolfKnauff2008RevisionStrategy}.

The reported mean conditional-choice percentages were 61.67 versus 22.08
for high versus low MP materials, with reported SDs 30.07 and 21.45;
for MT they were 82.50 versus 26.25, with SDs 20.13 and 24.86.
These are repeated-response summaries, not percentages of independent
people who always rejected the conditional. The paper's corresponding
test summaries do not supply recoverable t-statistics or usable exact
p-values, and no new confidence interval is reconstructed from the
aggregates. Its revision preferences are compatible with sensitivity
to credibility, but content selection prevents attributing everything
to an independently isolated probability parameter.

Experiment~2's report is explicitly preliminary: 40 people were planned,
but only 16 had been tested for this article, seven men and nine women,
aged 18--26 and monetarily compensated. It used 24 problem sets, four
items in each MP/MT by high/low/near-50\% condition, within participants.
Under the same uncertain-premise/certain-fact instruction, conditional
choice was 75.0 versus 37.5\% for high versus low MP and 81.3 versus
40.6\% for high versus low MT. The paper reports the paired contrasts
as \(z=-2.44,p=.015\) and \(z=-2.997,p=.003\), respectively.
Near-50\% items yielded 71.9\% conditional choice for MT versus 54.7\%
for MP, with \(z=-1.99,p=.047\).

The neutral MP result did \emph{not} show the predicted majority
preference for the categorical premise: conditional choice remained
54.7\%. The source describes only partial support for the mismatch
prediction. The review's discussion of a switch between strategies must
retain this exception and the available N=16, not present the planned
40 as the sample. Multiple content comparisons, no reported preregistered
analysis or contrast confidence intervals, and a single university
sample limit causal and population conclusions
\cite[pp.~834--836]{WolfKnauff2008RevisionStrategy}.
Decision time and preference alone do not directly measure two distinct
internal processing systems.

\subsection{Knauff and colleagues 2013, Experiment 1: what relocation meant}

This study tested 23 paid Giessen students individually, three male,
mean age 22.22 years with reported SD 2.45; two were psychology students.
The within-person task had 32 problems plus four practice trials, using
fruits and tools. Participants first read two spatial premises and selected
a compatible diagram, then judged a contrary new fact declared irrefutable,
then chose between two revised diagrams. Left/right relation forms and
object arrangements were counterbalanced. This was a constrained diagram
choice, not unrestricted navigation or spontaneous physical relocation.

Trials with an incorrect initial diagram or an incorrect contradiction
judgment were excluded from subsequent analysis. The source reports
95.65\% correct initial selections (SD 1.30) and 90.49\% correct contradiction
judgments (SD 9.01), but does not provide the surviving item-by-person
matrix. Among the analyzed revision responses, the LO-consistent option
had reported mean 86.67\% (SD 3.44) and the RO-consistent option 13.33\%
(SD 3.44), with \(z=-4.22,p<.001\) in the reported Wilcoxon test
\cite[Experiment~1]{KnauffEtAl2013SpatialRevision}.

LO and RO are the authors' located/reference roles in the premise
representation. The study supports a preference under those materials
and response choices. It does not establish that real-world agents always
move a located object, or uniquely determine a pictorial rather than
role-sensitive sentential format. An ordered predicate representation can
also distinguish its arguments if a cognitive account gives those roles
different processing policies. The chapter's separate discussion of size
and movability cites a later grounding study; its sample is not combined
with the 23 participants above
\cite[\S3.2]{GazzoCastanedaKnauff2021Defeasible}.

\section{Three explanatory programs, not three observed mental objects}
\label{sec:p19-r54-theories}

\paragraph{Mental models.}
The chapter distinguishes the three valuations allowed by a material
conditional from the initially explicit p-and-q model and an implicit
record of further possibilities. Retrieval can make a disabling case
explicit. Pragmatic modulation gives relevant background information
priority when the initial representation conflicts with it. This is a
proposal about interpretation and search, not a proof that a fixed
material conditional remains true in the new counterexample.

For revision the model account has three stages: search for a model
satisfying all the premises; identify an inconsistency if this fails;
revise using a mismatch between the contrary fact and the currently
explicit representation, then seek a diagnostic explanation. The source
uses this to contrast MP's new not-q with MT's new p. The prediction
depends on which model was represented, not just the external formulas.
The imperfect neutral-MP result above shows why agreement with one
direction of an asymmetry is not complete confirmation of that process
\cite[\S4.1]{GazzoCastanedaKnauff2021Defeasible}.

\paragraph{Probability.}
The probabilistic program models confidence in an ordinary conditional
by \(\Pr(q\mid p)\), and confidence in the reverse inference by
\(\Pr(p\mid q)\). The Ramsey-test description supposes the antecedent
and assesses the consequent. This conditional-probability identification
is a modeling commitment about ordinary conditionals, not an identity
with the probability of the material proposition \(p\to q\).
The chapter also emphasizes that frequency information can influence
reasoning without literal enumeration of exception representations.

\begin{hbtheorem}{thm:p19-r54-probability-boundary}{proved=annex:p19.r54.probability-boundary}{The source illustration, correct types and conditioning hypotheses}
For the chapter's illustrative weights 3,1,2,4 on
\((p,q)=(1,1),(1,0),(0,1),(0,0)\), respectively,
\(\Pr(q\mid p)=3/4\), \(\Pr(p\mid q)=3/5\),
\(\Pr(q)=1/2\) and \(\Pr(p\to q)=9/10\).
These are stipulated chances, not participant frequencies.
For a probability distribution with \(\Pr(\neg q)=1\),
\(\Pr(q\mid p)=0\) follows only when \(\Pr(p)>0\);
with zero denominator the ordinary ratio is undefined.
If \(\Pr(p)=1\), then \(\Pr(q\mid p)=\Pr(q)\),
not the event q itself. Ordinary conditioning on an incoming fact
also requires its prior probability to be positive.
\end{hbtheorem}

The chapter applies that calculation to possible MP--MT differences
after accepting a contrary fact with certainty. It is not yet a complete
psychological account of how the new distribution was obtained. In
particular, a categorical contradiction assigned prior probability zero
cannot be assimilated by ordinary conditionalization alone. A person
may instead change a model, a conditional confidence, a source assessment
or another assumption. Neither a forced relative-premise choice nor a
group endorsement proportion is automatically that person's measured
subjective probability \cite[\S4.2]{GazzoCastanedaKnauff2021Defeasible}.

\paragraph{Dual processes.}
The chapter's specified dual-process comparison associates heuristic
responses with overall likelihoods and more analytic responses with
retrieving particular disablers or alternatives. This is not a license
to rename every fast answer probabilistic and every slow answer logical.
Task format, time and working-memory demands are proposed moderators;
the cognitive mechanisms need independent discriminating evidence.
Wolf and Knauff's high/low versus neutral comparison is one relevant
test, with the partial results and sample boundary stated above
\cite[\S4.3]{GazzoCastanedaKnauff2021Defeasible}.

\section{Normative revision and the unification argument}
\label{sec:p19-r54-normative-bridge}

AGM in \cref{def:p08-agm} operates on a deductively closed classical
belief set, with the postulates in \cref{tab:p08-agm-postulates}.
Experimental sentence choices, graded doubts, diagnostic explanations
and cognitive representation changes are not automatically operations
on that same mathematical object. The laboratory instruction that an
incoming fact is unquestionable resembles a success constraint, but
does not demonstrate logical closure, extensionality or the whole AGM
postulate set. Entrenchment requires an ordering or selection policy;
the words ``minimal change'' alone do not select which premise to retain.

\begin{hbtheorem}{thm:p19-r54-two-revisions}{proved=annex:p19.r54.two-revisions}{Two opposite sentence choices can both satisfy the eight AGM conditions}
For the finite example, \(\operatorname{Cn}\) denotes classical closure
of a formula set and \(\operatorname{Th}(X)\) the formulas true at
every valuation in X; the empty X has the inconsistent theory.
In the finite language of p,q, fix \(K=\operatorname{Cn}(\{p\to q,p\})\),
whose only model is \((1,1)\). Define revision by selecting the
lowest-ranked input valuations and taking their theory. The total
orders
\[
 (1,1)\prec(1,0)\prec(0,0)\prec(0,1),\qquad
 (1,1)\prec(0,0)\prec(1,0)\prec(0,1)
\]
each yield all eight AGM conditions for this fixed K and arbitrary
formula inputs. On input \(\neg q\), the first retains p but rejects
\(p\to q\); the second retains \(p\to q\) but rejects p.
Thus a preference for either sentence alone does not prove or refute
the AGM conditions. Counting deletions in the two-sentence base gives
a one-deletion tie, and is not itself the AGM theory-level criterion.
\end{hbtheorem}

The spatial example supplies a similar contrast without asserting that
spatial relocation is an AGM theorem. From A--B--C and a new requirement
that A be right of C, either B--C--A or C--A--B retains one original
ordering statement and abandons the other. Both move one outer object.
A preference for the LO option therefore requires a further cognitive,
linguistic or practical policy beyond that equal movement/deletion count.

The chapter's unification argument contrasts forced premise rejection
with asking for an explanation or allowing graded doubt. It discusses
Dieussaert and colleagues' before/after conditional-believability ratings,
Walsh and Johnson-Laird's explanations, Khemlani and Johnson-Laird's
explanation choices, and Politzer and Carles's doubt option. These are
different response opportunities, not interchangeable demonstrations
that one belief-set operation was performed. A nonsignificant before/after
contrast is not an equivalence test and does not establish that no revision
occurred. Explaining a contradiction by an exception may change what a
rule is taken to mean, without amounting to complete rejection of that
rule. That is the source's proposal for a common role of disablers, not
a formal identity of defeasible consequence and AGM revision
\cite[\S5]{GazzoCastanedaKnauff2021Defeasible}.

\paragraph{Limits and research direction.}
The source closes by criticizing an exclusively computational account
that specifies an output quantity without specifying a cognitive process;
it also notes difficulties for the truth principle and objections to
unspecified dual-process divisions. Its call for collaboration among
psychology, philosophy and AI concerns explicit links between representations,
operations, tasks and observations. It does not infer an ought from a
response majority. An empirical test of a formal account needs a declared
mapping from task sentences and measurements to its mathematical states,
plus predictions not chosen only after seeing the response pattern.
Different populations, content, instructions and response formats require
their own evidence. The finite primary records above do not settle all
of the review's cited effects or generalize them to ordinary life as a whole
\cite[\S6]{GazzoCastanedaKnauff2021Defeasible}.
